\documentclass[11pt,letterpaper]{article}
\usepackage[margin=0.75in]{geometry}
\usepackage[utf8]{inputenc}
\usepackage[T1]{fontenc}
\usepackage{lmodern}
\usepackage{helvet}
\usepackage{textcomp}
\usepackage{parskip}
\usepackage{longtable}
\usepackage{booktabs}
\usepackage{array}
\usepackage{calc}
\usepackage{amssymb}
\usepackage{enumitem}
\usepackage{titlesec}
\usepackage{xcolor}
\usepackage[colorlinks=true,urlcolor=blue,linkcolor=black]{hyperref}
\renewcommand{\familydefault}{\sfdefault}
\setlist[itemize]{leftmargin=1.3em,itemsep=1pt,topsep=2pt}
\setlist[enumerate]{leftmargin=1.5em,itemsep=1pt,topsep=2pt}
\titleformat{\section}{\large\bfseries}{}{0em}{}[\vspace{0.1em}\hrule\vspace{0.2em}]
\newcolumntype{L}[1]{>{\raggedright\arraybackslash}p{#1}}
\providecommand{\tightlist}{\setlength{\itemsep}{0pt}\setlength{\parskip}{0pt}}
\setlength{\parindent}{0pt}\setlength{\parskip}{4pt}
\begin{document}
\section{Research Statement --- Ángel L.
Robles-Fernández}\label{research-statement-uxe1ngel-l.-robles-fernuxe1ndez}

\subsection*{Oregon State University --- Assistant Professor,
Integrative
Biology}\label{oregon-state-university-assistant-professor-integrative-biology}

\begin{center}\rule{0.5\linewidth}{0.5pt}\end{center}

My research program addresses a straightforward question: \emph{Can we predict which species will interact and where?} Answering this question requires integrating three dimensions of biodiversity: genomic and  phylogenetic information, environmental information, and species distribution modeling. My research seeks to integrate this information with computational methods by bringing together large volumes of data in open databases. My work has progressed along two complementary lines: first, the development of machine learning frameworks to predict interactions between hosts and pathogens on a continental and global scale; and, second, the creation of demographic models of species distribution that account for climate variability—a factor ignored by traditional approaches despite its proven importance for population viability.

At Oregon State University, I will bring these themes together in a research program that can be integrated with the mentoring of students of anylevel, using open-source software tools and available open data to train students in computational ecology, while also producing publishable scientific results.

\subsection*{Arc 1: Predicting species interactions from biodiversity
dimensions
(2017--2022)}\label{arc-1-predicting-species-interactions-from-biodiversity-dimensions-20172022}

The core insight driving my early work is that biotic interactions ---
who parasitizes whom, which hosts share pathogens --- are not random.
They are structured by evolutionary history, environmental overlap, and
geographic proximity. My contribution has been to formalize this insight into predictive frameworks.

My first paper on this topic (Robles-Fernández \& Lira-Noriega, 2017,
\emph{Frontiers in Applied Mathematics and Statistics}) established the
protocol: combine phylogenetic reconstructions of host plants with
geographic occurrence data to estimate a pest-host interaction index in geographic space. This work demonstrated that pest-host interactions are constrained by host phylogeny and that these constraints can be
projected onto maps to identify hotspots of vulnerability. The paper
also introduced \texttt{geotax}, an R package that operationalizes this
protocol for any pest-host system. We also have developed recently an updated protocol of this method base on PU learning, in a new R package called \texttt{phylopred}.

Building on this framework, my next study (Robles-Fernández,
Santiago-Alarcon \& Lira-Noriega, 2021, \emph{Frontiers in Veterinary
Science}) applied machine learning to predict dengue susceptibility in
wild mammals across the Americas as a function of phylogenetic,
environmental, and geographic distances. We found that bats were
predicted to be highly susceptible to DENV, consistent with field
studies, and that temperature was a dominant factor determining
transmission risk. Critically, we identified potential reservoir species
that had not been previously considered.

The generalizability of this approach became the focus of my
dissertation work, published in \emph{PNAS} (Robles-Fernández,
Santiago-Alarcon \& Lira-Noriega, 2022). Here, I tested the framework
across three fundamentally different host-pathogen systems: directly
transmitted coronaviruses in bats, and two vector-borne systems --- West
Nile Virus and avian malaria in birds. I showed that reliable global
susceptibility predictions could be obtained even when using less than
20\% of incidence data, and that different systems were driven by
different dimensions: coronavirus susceptibility was mostly geographic,
WNV mostly phylogenetic, and avian malaria mostly environmental. This
result has practical implications: it means that surveillance efforts
must be tailored to the specific ecology of each pathogen system, and
that a one-size-fits-all approach to disease forecasting will fail.

Two subsequent collaborations extended this framework to new systems. In
Krasnov et al.~(2025, \emph{Parasitology}), we analyzed determinants of
ectoparasite species richness (fleas and gamasid mites) on small mammals
across six biogeographic realms, finding that host trait similarity was
the strongest predictor of parasite sharing --- stronger than phylogeny,
geography, or environment. In Tinoco-Domínguez et al.~(2025,
\emph{American Journal of Botany}), we applied network analysis to
\emph{Phoradendron} mistletoes and their hosts, showing that the
probability of interaction depends primarily on host phylogenetic
distance. Across all these systems --- viruses, bacteria, ectoparasites,
and parasitic plants --- the same principle holds: biodiversity
dimensions structure biotic interactions, and this structure is
predictable.

\subsection*{Arc 2: Climate variability and species distributions
(2022--present)}\label{arc-2-climate-variability-and-species-distributions-2022present}

During my postdoctoral work in the Reuman Lab at the University of
Kansas, my focus shifted from predicting \emph{who interacts with whom}
to predicting \emph{where species can persist} under environmental
change. The shift was motivated by a recognized gap: traditional species
distribution models (SDMs) use long-term climate averages and ignore
interannual variability, despite the fact that ecological demographers
have long known that variability influences population viability.

In Hernández, Robles-Fernández \& Upham (2024, \emph{Ecography}; equally
contributing first authors), we addressed a related problem: how
historical environmental suitability shapes present-day genetic
diversity. We developed the Suitability Prevalence Area (SPA) method,
which projects Grinnellian niches back through the Late Quaternary to
identify areas of long-term environmental stability. Applying SPA to the
tassel-eared squirrel (\emph{Sciurus aberti}), we showed that
populations closer to the historical endemic area harbor higher genetic
diversity and lower genetic differentiation --- a result that connects
biogeographic history with population genomics and provides a predictive
framework for conservation prioritization.

The SPA work revealed that environmental \emph{stability} matters, but
it did not address the complementary question: how does environmental
\emph{variability} constrain distributions? To answer this, I
co-developed \textbf{XSDM} (Berti, Robles-Fernández et al., 2026,
\emph{bioRxiv}), a novel SDM framework that integrates stochastic
demography into species distribution modeling. XSDM maximizes a
likelihood function based on environmental time series and
presence/absence records to reconstruct species' fundamental niches and
project their potential geographic ranges. I am the creator and
maintainer of the \texttt{xsdm} R package (CRAN-ready, C++ backend with
RcppParallel), which implements this framework.

Our key results: climate variability reduces species' potential
distributions by an average of 22\%, and up to 45\% in some species.
Moreover, the sensitivity of distributions to projected changes in
climate \emph{variability} was comparable in magnitude to the
sensitivity to changes in climate \emph{means}. This has profound
implications: SDMs that ignore variability are systematically
overestimating species' ranges under both current and future conditions.
As global change alters not just mean temperatures but also the
frequency of extreme events, XSDM provides a tool to correct these
biases.

\subsection*{Future research at Oregon State University: Integrating field
ecology with global
synthesis}\label{future-research-at-appalachian-state-integrating-field-ecology-with-global-synthesis}

At Oregon State University, I will unify these two arcs into a research
program organized around a central question: \textbf{how does climate
variability reshape species interactions, and what are the consequences
for disease emergence risk?}

This program has three components, each designed to involve
undergraduate and master's students at every stage.

\subsubsection*{Component 1: Computational macroecology with public
data}\label{component-1-computational-macroecology-with-public-data}

My research program is based on publicly available data—GBIF records, GenBank sequences, xbioclim climate layers, and IUCN distribution maps—as well as my own mathematical and statistical models and the software I have already developed. For students at Appalachian State, this is an advantage. An undergraduate student with a laptop can learn basic R skills in my courses and how to download and handle occurrence records for North American mammals, extract climate data, and build a species distribution model using \texttt{xsdm} --- all within a single lab session. There is no barrier to entry. For their undergraduate thesis, students can expand this analysis to different clades or functional groups and conduct comparative studies—for example, by modeling different climate change scenarios and obtaining results that deepen our understanding of how climate variability affects species distribution. My line of research is global and computational in nature, making it infinitely scalable to accommodate undergraduate students' work.

\subsubsection*{Component 2: Mechanistic modeling with
XSDM}\label{component-2-mechanistic-modeling-with-xsdm}

The XSDM framework is ideally suited for projecting how climate
variability affects host-parasite systems. Traditional SDMs would model
a host species' distribution based on average climate; XSDM can model it
based on the full time series of climate variability, capturing the
demographic consequences of extreme years. I will apply XSDM to host and
vector species in the Southern Appalachians, projecting their
distributions under CMIP6 climate scenarios. Students will learn to run
SDMs, interpret output, and visualize uncertainty --- all using tools I
have already built and documented.

\subsubsection*{Component 3: Global synthesis of disease risk under
climate
change}\label{component-3-global-synthesis-of-disease-risk-under-climate-change}

The third component scales back up to the global level, connecting my
PNAS framework (predicting susceptibility from biodiversity dimensions)
with my XSDM framework (projecting distributions under variability). I
will ask: \emph{as climate variability increases, which regions will see
the greatest reshuffling of host communities, and where will novel
host-pathogen encounters become most likely?} This synthesis will use
publicly available data from GBIF, xbioclim, and the IUCN
database. All the data and software provided is accessible to undergraduates with basic programming skills. My developed R packages provide the computational infrastructure.

\subsection*{Student involvement and publication
strategy}\label{student-involvement-and-publication-strategy}

A computational research program in disease ecology is uniquely suited
for a primarily undergraduate institution. It requires no expensive
wet-lab infrastructure, no animal care protocols beyond standard field
methods, and no recurring consumable costs. Students can contribute from
their first semester: a freshman can clean occurrence data in R while a
senior runs XSDM projections for their capstone. I am committed to
mentoring undergraduates as co-authors, and I will formalize this by
establishing a lab structure where each student owns a component of the
research pipeline from data-cleaning to analysis to manuscript.

My target is 1--2 publications per year, with undergraduate and MS
students as co-authors. The projects are modular: a single semester of
trapping data can become a short communication in \emph{Biodiversity
Data Journal}; a two-semester SDM project can become a full paper in
\emph{Biodiversity informatics} or \emph{Diversity and Distributions}. I will pursue external funding through channels appropriate to my position --- including foundations, international and society-based schemes, and institutional internal support --- but the research program is designed to produce publishable results even without external grant support.


\subsection*{Integration with teaching}\label{integration-with-teaching}

This research program is integrated with my teaching curriculum. In Global Change Ecology, they will run XSDM projections and compare them to published SDM studies. In Evolution, they will learn the bases of phylogenetic comparative studies applied to host and parasite systems. Each course feeds directly into the research pipeline, and
the research pipeline feeds back into course materials. 

\subsection*{Summary}\label{summary}

My research program is based on publicly available open-source software and datasets (nine packages in R/C++ and web services in production), seven peer-reviewed publications (with other manuscripts under review), and a clear, long-term research trajectory: from predicting who interacts with whom, to predicting where species may persist, to projecting how climate change will reshape both of these aspects. At Oregon State University, I will build this program into a student-centered research group that produces rigorous science while training the next generation of computational ecologists.
\end{document}
